\documentclass[10pt,a4paper]{amsart}
\usepackage[margin=19mm]{geometry}
\usepackage{amsmath,amssymb,mathtools}
\usepackage[scr=rsfs,cal=euler]{mathalfa}
\usepackage{xfrac}
\usepackage[unicode,hidelinks,bookmarks=false]{hyperref}
\newcommand{\ie}{i.e.\ }
\newcommand{\cf}{cf.\ }
\newcommand{\resp}{respectively\ }
\newcommand{\Subsection}{\subsection}
\providecommand{\nobreakdash}{\nobreak\hbox{-}}
\setlength{\emergencystretch}{2em}
\multlinegap0pt
\usepackage{tikz}
\usepackage{amssymb}
\usepackage{enumerate}
\newtheorem{theorem}{Theorem}
\usepackage{cleveref}
\crefname{theorem}{Theorem}{Theorems}
\newcommand{\Z}{\mathbb{Z}}
\title{Counting Salem numbers arising from arithmetic hyperbolic orbifolds}
\author{Michelle Chu, Plinio G. P. Murillo, Otto Romero,, Lola Thompson}
\date{Source: arXiv:2508.08003v2 (CC BY 4.0)}
\begin{document}
\maketitle

\noindent\emph{Source and licence.} Counting Salem numbers arising from arithmetic hyperbolic orbifolds. Version arXiv:2508.08003v2 (CC BY 4.0), distributed by arXiv under the Creative Commons Attribution 4.0 International licence, http://creativecommons.org/licenses/by/4.0/, as recorded in the arXiv licence metadata for this exact version and shown on the abstract page https://arxiv.org/abs/2508.08003v2. Full licence notice: \texttt{licenses/arxiv-2508.08003-CC-BY-4.0.txt}.\par\smallskip

\noindent\textit{Editorial scope.} Counted unit: the article's theorem sol. (source0.tex lines 285--294). The article's complete original proof of it is delivered with it (source0.tex lines 357--433, one \texttt{proof} environment of 77 lines, which the article places away from the statement). No mathematics is written, repaired or re-derived: the statement and the proof are the article's own lines, in the article's own notation, and no retained line is substituted or re-flowed.  No further source-local context is needed: every cross-reference of every retained line resolves inside the two delivered spans.  Nothing was omitted from any retained span, so the package carries no omission record.  No retained line contains a citation, so the wrapper carries no bibliography and no bibliography entry.  No retained line contains a figure environment, an image-inclusion command or drawing code, so no image is delivered and the package declares no asset dependency.  Editorial formatting only, in addition: an AMS \texttt{amsart} preamble that restates the article's own declarations and macros used by the retained lines, namely the environments \texttt{theorem} on separate counters, together with the standard packages those lines need.  The retained environments print in this document as Theorem 1 in order of appearance, whereas the article prints the same items with the numbering of its own full document.  The complete native source of the article as distributed in the arXiv e-print for this exact version is bundled unchanged as \texttt{sources/183064/source0.tex}.
\begin{theorem}\label{th:diop. sol.}
    Let $D$ be a square-free positive integer and $t$ be the number of distinct prime factors of $D$. Denote the number of integer solutions to $A^2+DB^2=C^2$ with $\gcd(A,C)=1$ and $|C|\leq X$ by
    \[ P_D(X)=\#\{(A,B,C)\in \Z^{3}\mid A^2+DB^2=C^2, \, \gcd(A,C)=1,  \text{  and  } |C|\leq X\}.
    \]
    Then
    \[ P_D(X) = \begin{cases}
    2^{t}\frac{8X}{\pi\sqrt{D}} + O(\sqrt{X}\log X)   &\text{ if $D$ is odd},\\
    2^{t}\frac{6X}{\pi\sqrt{D}} + O(\sqrt{X}\log X)   &\text{ if $D$ is even} .
    \end{cases} \]
\end{theorem}

\begin{proof}[Proof of \Cref{th:diop. sol.}]
Let $D$ be positive and square-free. We will count integer solutions to the Diophantine equation $A^2 + DB^2 = C^2$ where $\gcd(A, C) = 1$ and $C \leq X$. We begin with the classical method of drawing an ellipse $A^2 + DB^2 = 1$ and drawing a reference point along the ellipse (for example, $(1, 0)$), and then drawing a line from this reference point to a rational point $r$ along the $y$-axis. The point where the line intersects the ellipse will also be a rational solution to $A^2 + DB^2 = 1$. A line through both $(1,0)$ and $r$ is given by the equation $B = (1-A)r$. By substituting $B = (1-A)r$ into the Diophantine equation and solving for $A$ and $B$, we obtain $A = \frac{-1 + Dr^2}{1 + Dr^2}$ and $B = \frac{2r}{1+Dr^2}$. Setting $r=V/U$ yields the parametrization 
    \begin{equation} \label{eq:all sols param}
    A= S\frac{-U^2+DV^2}{T},\quad B=S\frac{2UV}{T}, \text{ and } C=S\frac{U^2+DV^2}{T}, \end{equation}
where $S,U,V$ are integers satisfying $\gcd(U,V)=1$, and $T$ is the greatest common divisor of the numerators.

Suppose $(A,B,C)$ is a positive primitive solution written as in \eqref{eq:all sols param}. We may then assume that $S=1$, $U,V\geq0$, and $U^2\leq DV^2$. Let $D_1=\gcd\{U^2+DV^2,UV\}$. Since $\gcd(U,V)=1$, it must be that $D_1$ divides both $D$ and $U$. There is some $D_2>0$ and $u\geq0$ such that $D=D_1D_2$ and $U=D_1u$. It follows that any positive primitive solution can be written as 
\begin{equation} \label{eq:pos prim param}
     (A,B,C)=\frac{1}{\tau}(-D_1u^2+D_2v^2,2uv,D_1u^2+D_2v^2) 
\end{equation}
where $D_1 > 0$, $D_2 >0$, $D=D_1D_2$, $u \geq 0$, $v \geq 0$, $\gcd(u, v) = 1$, $D_1u^2\leq D_2v^2$, and $\tau=1$ when $B$ is even or otherwise $\tau=2$ and $u,v$ are odd when $B$ is odd.

We will now count the number of positive primitive solutions. We will consider four cases, depending on the parity of $D$ and on the relationship between $A$ and $C$. Notice that, if $\gcd(A, C) = 1$, then $\gcd((C-A), (C+A)) = 1$ or $2$. 

\vspace{10pt}
\noindent \emph{Case I:} $D$ is even and $A\equiv C\bmod2$. 

Then $B$ must be even. 
By \eqref{eq:pos prim param}, there are unique integers $D_1 > 0, D_2 >0$, $u \geq 0$, and $v \geq 0$ with $\gcd(u, v) = 1$ for which 
\[ A = D_2 v^2 - D_1 u^2, \hspace{0.2 in} B = 2uv, \hspace{0.2 in} C = D_1u^2 + D_2v^2.\]
We will show that every choice of such $D_1, D_2, u, v$ produces a positive primitive solution.

The congruence condition on $A$ and $C$ implies that $\gcd(C-A, C+A) = 2$. Therefore, there is no odd prime that divides both $A$ and $C$. If $A$ and $C$ are both even then $D_1 u^2$ and $D_2 v^2$ are both odd, which contradicts the fact that $D = D_1 D_2$ is even. As a result, for each choice of $D_1$ and $D_2$, we will count the elements in the set 
\[ \{(u, v) \in (\mathbb{Z}_{\geq 0})^2 : \gcd(u, v) = 1, D_1u^2 + D_2v^2 \leq X, D_1u^2 \leq D_2 v^2\}.\]
This count is precisely $N'(\sqrt{X}, \mathcal{C})$ where 
\[ \mathcal{C} = \{(x,y) \in (\mathbb{R}_{\geq 0})^2 : D_1x^2 + D_2y^2 \leq 1, D_1x^2 \leq D_2 y^2\}. \]
Observe that 
\[ \mathrm{vol}(\mathcal{C}) = \frac{\pi}{8 \sqrt{D_1 D_2}} = \frac{\pi}{8 \sqrt{D}}. \]
From our formula for $N'(X, \mathcal{C})$, we obtain 
\[ N'(\sqrt{X}, \mathcal{C}) \sim \frac{6}{\pi^2} \times \frac{\pi}{8\sqrt{D}} \times (\sqrt{X})^2 \sim \frac{3X}{4 \pi \sqrt{D}}. \]
As we run over all choices of $D_1$, our count becomes 
\[ 2^t \frac{3X}{4 \pi\sqrt{D}} + O(\sqrt{X} \log X), \]
where the error term comes from the formula for $N'(\sqrt{X}, \mathcal{C})$.

\vspace{10pt}
\noindent \emph{Case II:} $D$ is even and $A\not\equiv C\bmod2$. 

This case does not occur, since it would contradict $C^2-A^2=D B^2$.

\vspace{10pt}
\noindent \emph{Case III:} $D$ is odd and $A\equiv C\bmod2$.

Since $A \equiv C\bmod2$, then $B$ must be even. By \eqref{eq:pos prim param}, there are unique integers $D_1 > 0, D_2 >0$, $u \geq 0$, and $v \geq 0$ with $\gcd(u, v) = 1$ for which 
\[ A = D_2 v^2 - D_1 u^2, \hspace{0.2 in} B = 2uv, \hspace{0.2 in} C = D_1u^2 + D_2v^2.\]

Each choice of such $D_1, D_2, u, v$ produces a solution to $A^2+DB^2=C^2$. However, not all choices produce a primitive solution. The greatest common divisor may be $2$, which occurs if both $u$ and $v$ are odd. 

Thus, for each choice of $D_1, D_2$, we count the elements in the set 
\[ \{(u,v) \in (\mathbb{Z}_{\geq 0})^2 \mid \gcd(u, v) = 1, D_1u^2 + D_2v^2 \leq X, D_1u^2 \leq D_2v^2\} \]
and subtract off the number of elements in the set 
\[ \{(u,v) \in (2\mathbb{Z}_{\geq 0} + 1)^2 \mid \gcd(u, v) = 1, D_1u^2 + D_2v^2 \leq X, D_1u^2 \leq D_2v^2\}.\]
This is precisely given by 
\[ N'(\sqrt{X}, \mathcal{C}) - M'(\sqrt{X}, \mathcal{C}),\]
where $\mathcal{C}$ is the same as in \textit{Case I}. Using our formulae for $N'$ and $M'$, we obtain 
\[ N'(\sqrt{X}, \mathcal{C}) \sim \frac{3X}{4 \pi \sqrt{D}},\]
as in \textit{Case I}, minus 
\[ M'(\sqrt{X}, \mathcal{C}) \sim \frac{2}{\pi^2} \times \frac{\pi}{8 \sqrt{D}} \times (\sqrt{X})^2 \sim \frac{X}{4 \pi\sqrt{D}}. \]
This difference is asymptotically $\frac{X}{2 \pi \sqrt{D}}$. As we run over all choices of $D_1$, we obtain a contribution of $2^t \frac{X}{2\pi \sqrt{D}} + O(\sqrt{X} \log X).$\\

\vspace{10pt}
\noindent \emph{Case IV:} $D$ is odd and $A\not\equiv C\bmod2$.

Then $B$ must be odd and \(\gcd(C+A, C-A) = 1\). By \eqref{eq:pos prim param}, there are unique integers $D_1 > 0, D_2 >0$, $u \geq 0$, and $v \geq 0$ with $\gcd(u, v) = 1$ and both $u,v$ odd for which 
\[ A = \frac{1}{2}(D_2 v^2 - D_1 u^2), \hspace{0.2 in} B = uv, \hspace{0.2 in} C = \frac{1}{2}(D_1u^2 + D_2v^2).\]
Every choice of such $D_1, D_2, u, v$ produces a positive primitive solution to $A^2+DB^2=C^2$. Indeed, if an odd prime $p$ is a common divisor of $A$ and $C$, then $p$ also divides $A\pm C$, which is not possible since $\gcd(C+A, C-A) = 1$.

We also have
\[ 2C = D_1 u^2 + D_2 v^2 \leq 2X.\]
Then, for each choice of $D_1, D_2$, we count the elements in the set 
\[ \{(u, v)\in (2\mathbb{Z}_{\geq 0} + 1)^2: \gcd(u, v)=1, D_1u^2 + D_2v^2 \leq 2X, D_1u^2 \leq D_2v^2\}.\]
This is precisely 
\[ M'(\sqrt{2X}, \mathcal{C}) \sim \frac{2}{\pi^2} \times \frac{\pi}{8\sqrt{D}} \times (\sqrt{2X})^2 \sim \frac{X}{2 \pi \sqrt{D}}.\]
As we run over all choices of $D_1$, we obtain a total count of $2^t \frac{X}{2\pi\sqrt{D}} + O(\sqrt{X} \log X).$

\vspace{10pt}
Collecting the information from all four cases, we see that the contribution from even values of $D$ yields $2^t \frac{3X}{4\pi \sqrt{D}} + O(\sqrt{X} \log X)$ solutions to the Diophantine equation $A^2 + DB^2 = C^2$. On the other hand, the contribution from odd values of $D$ yields $2^t \frac{X}{\pi \sqrt{D}} + O(\sqrt{X} \log X)$ solutions. Observe that the above arguments only count positive solutions. Allowing for all $2^3$ combinations of $\pm$ solutions yields the result in \Cref{th:diop. sol.}.
\end{proof}

\end{document}
